\documentclass[12pt, a4paper]{article}

\usepackage[margin=2.5cm]{geometry}
\usepackage{amsmath, amssymb}
\usepackage{hyperref}
\usepackage{enumitem}
\usepackage{booktabs}
\usepackage{xcolor}
\usepackage{listings}
\usepackage{graphicx}
\usepackage{fancyhdr}
\usepackage{parskip}

\hypersetup{
    colorlinks=true,
    linkcolor=blue,
    urlcolor=blue,
    citecolor=blue
}

\lstset{
    basicstyle=\ttfamily\small,
    backgroundcolor=\color{gray!8},
    frame=single,
    breaklines=true,
    showstringspaces=false,
    keywordstyle=\color{blue},
    commentstyle=\color{green!50!black},
    stringstyle=\color{orange!80!black},
}

\pagestyle{fancy}
\fancyhf{}
\lhead{iBot Mini Project 1}
\rhead{Week 1: Foundations and Baseline}
\cfoot{\thepage}

% ------------------------------------------------------------------
\begin{document}

\begin{center}
    {\LARGE \textbf{Week 1: Foundations and Baseline}}\\[0.4em]
    {\large iBot Mini Project 1 -- Train a Robot to Walk}\\[0.2em]
    {\normalsize Estimated time: 4 to 5 hours}
\end{center}

\hrule
\vspace{0.5em}

\tableofcontents
\newpage

% ==================================================================
\section{Goals for This Week}

By the end of Week 1 you should be able to:
\begin{itemize}
    \item Explain what reinforcement learning is and how it is different from supervised learning.
    \item Define a Markov Decision Process (MDP) and identify its components in our walking robot task.
    \item Describe what a policy and a value function are.
    \item Set up the project environment and run the Walker2D simulation.
    \item Train a baseline PPO agent and view its training curve in TensorBoard.
    \item Record a short video of your trained agent.
\end{itemize}

% ==================================================================
\section{What is Reinforcement Learning?}

\subsection{The Core Idea}

In supervised learning you are given labelled examples and you teach a model to map inputs to outputs. In reinforcement learning (RL) there are no labels. Instead, an \textbf{agent} interacts with an \textbf{environment} by taking \textbf{actions}. After each action the environment gives the agent a \textbf{reward}: a single number that says how good or bad the action was.

The agent has one goal: collect as much reward as possible over time.

\medskip
\begin{center}
\begin{tabular}{ccc}
    \textbf{Agent sends:} & $\longrightarrow$ & action $a_t$ \\
    \textbf{Environment returns:} & $\longrightarrow$ & new state $s_{t+1}$, reward $r_t$
\end{tabular}
\end{center}
\medskip

This loop repeats. The agent gradually learns which actions lead to more reward.

\subsection{Why is RL Useful for Walking?}

You cannot write down an explicit program for walking because walking depends on the exact state of every joint and the terrain at every moment. Instead, you define a reward for walking forward and let the agent figure out the rest. The same approach works for:
\begin{itemize}
    \item Real robots (Boston Dynamics, ETH Zurich quadrupeds)
    \item Game characters
    \item Robotic arms, drones, and prosthetics
\end{itemize}

\subsection{Video Resource}

\textbf{Recommended:} Watch the first 20 minutes of the following lecture before reading further.

\begin{itemize}
    \item David Silver's RL Lecture 1 (Introduction to RL), available on YouTube at:\\
          \url{https://www.youtube.com/watch?v=2pWv7GOvuf0}\\
          This lecture gives clear, visual explanations of the agent-environment loop.
\end{itemize}

% ==================================================================
\section{Markov Decision Processes}

\subsection{Definition}

Every RL problem can be described as a \textbf{Markov Decision Process (MDP)}. Formally:

\begin{center}
\begin{tabular}{ll}
    \toprule
    Symbol & Meaning \\
    \midrule
    $\mathcal{S}$ & State space: the set of all possible states of the environment \\
    $\mathcal{A}$ & Action space: the set of all possible actions the agent can take \\
    $P(s' | s, a)$ & Transition function: probability of moving to state $s'$ from $s$ after action $a$ \\
    $R(s, a, s')$ & Reward function: the reward received after that transition \\
    $\gamma \in [0,1)$ & Discount factor: how much the agent values future rewards \\
    \bottomrule
\end{tabular}
\end{center}

\textbf{Markov property:} the next state depends only on the current state and action, not on the full history. This is why the state description must be complete: our 22-dimensional observation captures everything the agent needs.

\subsection{Our Walking Robot as an MDP}

\begin{center}
\begin{tabular}{lp{9cm}}
    \toprule
    MDP Component & Our Project \\
    \midrule
    $\mathcal{S} \subset \mathbb{R}^{22}$ & Joint angles, joint velocities, torso height, forward velocity, foot contacts, and more \\
    $\mathcal{A} \subset [-1,1]^6$ & Torques applied to 6 joints (scaled to $\pm 100$ Nm internally) \\
    $P$ & Implemented by the PyBullet physics engine (simulates gravity, collisions, joint dynamics) \\
    $R$ & Our reward function, which we design in Week 2 \\
    $\gamma = 0.99$ & Standard choice for locomotion tasks \\
    \bottomrule
\end{tabular}
\end{center}

\subsection{The Discount Factor}

The discount factor $\gamma$ controls how far into the future the agent cares. The \textbf{discounted return} from time step $t$ is:
\[
G_t = \sum_{k=0}^{\infty} \gamma^k r_{t+k+1}
\]

With $\gamma = 0.99$, a reward 100 steps in the future is worth $0.99^{100} \approx 0.37$ of an immediate reward. Choosing $\gamma$ matters:

\begin{itemize}
    \item $\gamma$ too low (e.g.\ 0.9): the agent focuses only on immediate rewards and may stand still rather than learn to walk.
    \item $\gamma$ too high (e.g.\ 0.999): small errors in value estimation compound, making training unstable.
    \item $\gamma = 0.99$: the standard for locomotion and most continuous control tasks.
\end{itemize}

% ==================================================================
\section{Policies and Value Functions}

\subsection{The Policy}

A \textbf{policy} $\pi_\theta(a \mid s)$ is the agent's strategy. It is a probability distribution over actions given the current state. In deep RL the policy is a neural network with parameters $\theta$.

During training the policy is \textbf{stochastic}: it samples an action from the distribution $\pi_\theta(\cdot \mid s)$. This randomness is essential for exploration.

During evaluation the policy can be made \textbf{deterministic}: take the action with the highest probability, $a = \arg\max_a \pi_\theta(a \mid s)$.

\subsection{The Value Function}

The \textbf{state value function} $V^\pi(s)$ is the expected total discounted reward the agent will collect from state $s$, if it follows policy $\pi$:
\[
V^\pi(s) = \mathbb{E}_\pi \left[ G_t \mid s_t = s \right]
\]

PPO trains both a policy network and a \textbf{value network} $V_\phi(s)$. The value network is used to estimate advantages (see Section~\ref{sec:ppo}).

\subsection{The RL Objective}

The goal of RL is to find parameters $\theta^*$ that maximise the expected return:
\[
J(\theta) = \mathbb{E}_{\tau \sim \pi_\theta} \left[ \sum_{t=0}^{T} \gamma^t r_t \right]
\]
where $\tau = (s_0, a_0, r_0, s_1, a_1, r_1, \ldots)$ is a trajectory (a full episode).

% ==================================================================
\section{Proximal Policy Optimization}\label{sec:ppo}

\subsection{The Problem with Simple Policy Gradients}

The most straightforward way to maximise $J(\theta)$ is gradient ascent:
\[
\theta \leftarrow \theta + \alpha \nabla_\theta J(\theta)
\]

The gradient can be estimated from experience. However, this \textbf{vanilla policy gradient} (also called REINFORCE) has a major weakness: a large gradient step can completely destroy a good policy. If you take one bad update, the policy may become so poor that it never recovers.

\subsection{The PPO Fix: Clipped Objective}

PPO fixes this by \textbf{clipping} the update. Define the probability ratio:
\[
r_t(\theta) = \frac{\pi_\theta(a_t \mid s_t)}{\pi_{\theta_{\text{old}}}(a_t \mid s_t)}
\]

This ratio tells you how much the new policy differs from the old policy on the action that was taken.

The \textbf{advantage} $\hat{A}_t$ estimates how much better action $a_t$ was compared to the average action at state $s_t$:
\[
\hat{A}_t = G_t - V_\phi(s_t)
\]

PPO's clipped objective is:
\[
\mathcal{L}^{\text{CLIP}}(\theta) = \mathbb{E}_t \left[ \min\!\left( r_t(\theta) \hat{A}_t,\; \text{clip}(r_t(\theta), 1 - \varepsilon, 1 + \varepsilon) \hat{A}_t \right) \right]
\]

where $\varepsilon = 0.2$. The clip prevents the ratio from going too far from 1, which keeps updates \textbf{proximal} (close) to the old policy. This is where PPO gets its name.

\subsection{Generalized Advantage Estimation}

Instead of using raw returns $G_t$, PPO uses \textbf{Generalized Advantage Estimation (GAE)} to reduce variance:
\[
\hat{A}_t^{\text{GAE}(\gamma, \lambda)} = \sum_{l=0}^{\infty} (\gamma \lambda)^l \delta_{t+l}
\qquad \text{where} \quad
\delta_t = r_t + \gamma V(s_{t+1}) - V(s_t)
\]

The parameter $\lambda$ trades off variance against bias:
\begin{itemize}
    \item $\lambda = 0$: uses only the one-step TD error (high bias, low variance).
    \item $\lambda = 1$: uses the full Monte Carlo return (low bias, high variance).
    \item $\lambda = 0.95$: the recommended balance.
\end{itemize}

\textbf{Paper (optional reading):}
\begin{itemize}
    \item PPO paper: Schulman et al.\ (2017), \textit{Proximal Policy Optimization Algorithms}. \url{https://arxiv.org/abs/1707.06347}
    \item GAE paper: Schulman et al.\ (2015), \textit{High-Dimensional Continuous Control Using Generalized Advantage Estimation}. \url{https://arxiv.org/abs/1506.02438}
\end{itemize}

% ==================================================================
\section{The Walker2D Environment}

\subsection{What is PyBullet?}

PyBullet is a free, open-source physics engine for Python. It simulates rigid body dynamics, gravity, joint torques, and collisions. We use it to simulate the bipedal robot. It runs entirely on CPU so no GPU is needed.

\subsection{The Robot}

Our robot is a 2D bipedal walker. It has:
\begin{itemize}
    \item A rigid torso body.
    \item Two legs, each with three joints: hip, knee, and ankle (6 joints total).
    \item Gravity set to $-9.81$ m/s$^2$.
    \item Physics running at 240 Hz; the agent sends a new action every 4 physics steps (60 Hz).
\end{itemize}

\subsection{Observation Space (22 dimensions)}

\begin{center}
\begin{tabular}{cll}
    \toprule
    Index & Feature & Description \\
    \midrule
    0 & Torso height & $z$-position of the torso centre of mass (m) \\
    1 & Torso pitch & Forward/backward tilt angle (rad) \\
    2--7 & Joint angles & Angle of each of the 6 joints (rad) \\
    8--13 & Joint velocities & Angular velocity of each joint (rad/s) \\
    14 & Forward velocity $v_x$ & How fast the robot moves forward (m/s) \\
    15 & Lateral velocity $v_y$ & Sideways movement (m/s) \\
    16 & Vertical velocity $v_z$ & Up/down movement (m/s) \\
    17 & Pitch rate & Torso angular velocity (rad/s) \\
    18 & Left foot contact & Binary: is the left foot touching the ground? \\
    19 & Right foot contact & Binary: is the right foot touching the ground? \\
    20--21 & Thigh angles & Repeated joint angles for richer signal \\
    \bottomrule
\end{tabular}
\end{center}

\textbf{Why normalization matters.} Raw observations span very different scales: heights are around 1 m, joint velocities can reach 5 rad/s, and contacts are binary (0 or 1). Without normalisation the neural network struggles with mixed-scale inputs. We use \texttt{VecNormalize} from Stable-Baselines3, which computes a running mean and standard deviation and normalises observations automatically.

\subsection{Action Space (6 dimensions)}

The action is a 6-dimensional vector of joint torques in $[-1, 1]$, scaled internally to $[-100, +100]$ Nm. During training, actions are sampled from a Gaussian policy $\pi_\theta(a \mid s) = \mathcal{N}(\mu_\theta(s), \sigma^2)$. During evaluation, the deterministic action $\mu_\theta(s)$ is used.

\subsection{Episode Termination}

An episode ends when:
\begin{enumerate}
    \item The torso height drops below 0.5 m (the robot fell over).
    \item The torso pitch exceeds $|1.2|$ radians (approximately 69 degrees from vertical, the robot tipped).
    \item 1000 steps have elapsed (the episode is truncated).
\end{enumerate}

% ==================================================================
\section{Default Reward Function}

In Week 2 you will design your own reward functions. For now, we use the \textbf{dense reward}:

\[
r_t = \underbrace{c_1 \cdot \text{clip}(v_x, 0, 5)}_{\text{forward velocity}}
    + \underbrace{c_2}_{\text{alive bonus}}
    - \underbrace{c_3 \|a\|_2^2}_{\text{energy penalty}}
    - \underbrace{c_4 |v_y|}_{\text{lateral penalty}}
    + \underbrace{c_5 \cdot \max(0, 1 - |h - h^*|)}_{\text{height bonus}}
\]

with defaults $c_1 = 1.0$, $c_2 = 1.0$, $c_3 = 0.001$, $c_4 = 0.5$, $c_5 = 0.1$, $h^* = 1.2$ m.

Each term has a specific purpose:
\begin{itemize}
    \item \textbf{Forward velocity reward:} the main signal. Clipped at 5 m/s to prevent the agent from optimising for unrealistic speeds.
    \item \textbf{Alive bonus:} staying upright earns +1 every step. Without this, a struggling agent might decide to fall quickly.
    \item \textbf{Energy penalty:} small penalty on large torques. This leads to smoother, more natural motion.
    \item \textbf{Lateral penalty:} discourages sideways drift. The robot should walk in a straight line.
    \item \textbf{Height bonus:} softly encourages maintaining proper posture.
\end{itemize}

% ==================================================================
\section{The Technology Stack}

\begin{center}
\begin{tabular}{lll}
    \toprule
    Component & Tool & Purpose \\
    \midrule
    Physics simulation & PyBullet 3.x & Simulate robot dynamics \\
    RL algorithm & PPO via Stable-Baselines3 & Train the policy \\
    Environment API & Gymnasium 0.29 & Standard interface for RL environments \\
    Logging & TensorBoard & Track and plot training metrics \\
    Plotting & Matplotlib & Generate report figures \\
    Video & imageio + ffmpeg & Record MP4 demos \\
    Compute & CPU only & No GPU needed \\
    \bottomrule
\end{tabular}
\end{center}

% ==================================================================
\section{Setting Up the Project}

\subsection{Installation}

\begin{lstlisting}[language=bash]
# Create a virtual environment
python -m venv venv
source venv/bin/activate    # Windows: venv\Scripts\activate

# Install all packages
pip install -r requirements.txt
\end{lstlisting}

\subsection{Running a Quick Test}

Run this from the project root to check that the environment loads:

\begin{lstlisting}[language=python]
from env.walker_env import Walker2DEnv
import numpy as np

env = Walker2DEnv()
obs, info = env.reset()
print("Observation shape:", obs.shape)   # should be (22,)
print("First few values:", obs[:5])

for _ in range(10):
    action = env.action_space.sample()   # random action
    obs, reward, terminated, truncated, info = env.step(action)
    print(f"reward={reward:.3f}  height={obs[0]:.3f}")
    if terminated:
        break

env.close()
print("Environment test passed.")
\end{lstlisting}

% ==================================================================
\section{Understanding TensorBoard}

After you start training (see Section~\ref{sec:baseline}), open TensorBoard to watch the agent improve.

\begin{lstlisting}[language=bash]
tensorboard --logdir runs/
\end{lstlisting}

Key metrics to watch:

\begin{center}
\begin{tabular}{lll}
    \toprule
    Metric & TensorBoard Tag & What to look for \\
    \midrule
    Episode reward & \texttt{rollout/ep\_rew\_mean} & Should rise over time \\
    Episode length & \texttt{rollout/ep\_len\_mean} & Should increase as robot survives longer \\
    Policy loss & \texttt{train/policy\_loss} & Should oscillate slightly but not diverge \\
    Value loss & \texttt{train/value\_loss} & May rise initially as the value function learns harder states \\
    Entropy & \texttt{train/entropy\_loss} & Measures exploration level; should gradually decrease \\
    Approx KL & \texttt{train/approx\_kl} & Keep this below 0.02; large values mean unstable updates \\
    Clip fraction & \texttt{train/clip\_fraction} & Fraction of clipped gradients (ideally below 0.2) \\
    Explained variance & \texttt{train/explained\_variance} & How well the value function predicts returns; should rise toward 1 \\
    \bottomrule
\end{tabular}
\end{center}

\textbf{Normal vs.\ worried:}
\begin{itemize}
    \item Reward rising slowly for the first 200k steps: this is normal.
    \item Reward flat for 500k+ steps: likely a reward function or hyperparameter issue.
    \item Reward rising then suddenly collapsing: the learning rate may be too high.
\end{itemize}

% ==================================================================
\section{Your Task: Train the Baseline Agent}\label{sec:baseline}

\begin{enumerate}
    \item Open \texttt{starter\_week1.ipynb} in Jupyter.
    \item Follow the cells to set up and run your first training session.
    \item Train for at least 500\,000 steps (1\,000\,000 preferred for clear results).
    \item Take a screenshot of the TensorBoard \texttt{ep\_rew\_mean} curve.
    \item Record a short video of your agent.
    \item Answer the reflection questions at the end of the notebook.
\end{enumerate}

As an alternative to the notebook, you can run the command-line script:
\begin{lstlisting}[language=bash]
python train.py --reward dense --total_steps 1000000
\end{lstlisting}

% ==================================================================
\section{Deliverables for Week 1}

\begin{enumerate}[label=\textbf{D\arabic*.}]
    \item \textbf{Environment running:} the test cell in \texttt{starter\_week1.ipynb} runs without errors.
    \item \textbf{Training started:} you have a run folder under \texttt{runs/} with TensorBoard logs.
    \item \textbf{TensorBoard screenshot:} include this in your weekly notes or report.
    \item \textbf{Baseline reward recorded:} note the mean episode reward at 500k steps.
    \item \textbf{Video (optional but recommended):} record 10 to 20 seconds of your agent from \texttt{evaluate.py}.
    \item \textbf{Reflection answers:} fill in the written questions in the starter notebook.
\end{enumerate}

% ==================================================================
\section{Further Reading}

These are not required, but they are useful if you want to understand the material more deeply.

\begin{itemize}
    \item \textbf{Sutton and Barto (2018).} \textit{Reinforcement Learning: An Introduction.} Free PDF at \url{http://incompleteideas.net/book/the-book.html}. Read Chapters 1 to 4 for the MDP and value function foundations.

    \item \textbf{Spinning Up (OpenAI).} A practical RL guide with code. \url{https://spinningup.openai.com/en/latest/}

    \item \textbf{Stable-Baselines3 documentation.} \url{https://stable-baselines3.readthedocs.io}

    \item \textbf{YouTube:} ``Reinforcement Learning in 3 Hours'' by Nicholas Renotte covers training CartPole and more, at a beginner level. Search for it on YouTube.

    \item \textbf{YouTube:} The series ``RL Course by David Silver'' (DeepMind, available on YouTube) is the best free lecture series on RL. Lectures 1 and 2 cover all the math in Sections 3 and 4 of this guide.
\end{itemize}

\end{document}
